\keys{multiobjective optimization, SPEA2, in vitro fertilization, memetic algorithm, fitness landscape analysis}

\begin{abstract}{IVF/SPEA2: A Hybrid SPEA2 Algorithm with In Vitro Fertilization for Multi-Objective Optimization}
SPEA2 assigns fitness by combining accumulated dominance strength with a density term based on the distance to the $k$-th nearest neighbour. This mechanism gives the algorithm a well-distributed solution set, but it penalises exactly those offspring that local search produces, because they are born close to their parents. This dissertation proposes IVF/SPEA2, a coupling of the in vitro fertilization method to SPEA2 that addresses this difficulty through two host-specific decisions: assigning each mother a distinct father chosen by distance in objective space, and making further intensification cycles conditional on an improvement in the population's mean fitness rather than on the quality of a single offspring. An activation gate based on cumulative budget keeps the number of evaluations per generation unchanged, so that the comparison is evaluation-normalised. The evaluation covers 51 instances from the ZDT, DTLZ, WFG and MaF suites, with 60 runs per configuration and a budget of 100,000 evaluations, reporting inverted generational distance as the primary endpoint and hypervolume as a mandatory secondary one, under Holm multiplicity correction. Against canonical SPEA2, IVF/SPEA2 obtains 22 wins, 3 ties and 3 losses on the bi-objective instances and 15 wins, 7 ties and 1 loss on the tri-objective ones; restricted to the 39 instances not used for parameter calibration, the advantage holds at 19 and 12 wins respectively. Gains concentrate on regular front geometries and weaken or reverse on disconnected ones. Complementary studies bound this result: the benefit depends on the host algorithm, being far smaller when the same intensification scheme is coupled to NSGA-II; transfer to three constrained engineering problems yields mixed outcomes; and the ablation does not separate the full coupling from the variant omitting both decisions, so their justification remains mechanistic. An activation analysis shows that static landscape descriptors rank the expected magnitude of the gain but do not support a binary decision that generalises across problem families, whereas archive turnover measured in the early generations supports a threshold rule with balanced accuracy of 0.754. The work concludes that the coupling is a bounded, geometry-dependent improvement rather than a universal enhancement of SPEA2.
\end{abstract}
